\documentclass[11pt]{amsart}
\usepackage{amssymb}
\usepackage[margin=1.15in]{geometry}
\usepackage{tabularx}
\usepackage[hidelinks]{hyperref}
\hypersetup{pdftitle={Finite projective modules and flatness}}

\newtheorem{theorem}{Theorem}[section]
\newtheorem{lemma}[theorem]{Lemma}
\theoremstyle{definition}
\newtheorem{definition}[theorem]{Definition}
\newcommand{\id}{\operatorname{id}}
\newcommand{\Span}{\operatorname{Span}}
\newcommand{\code}[1]{\texttt{#1}}

\title{Finite projective modules and flatness}
\date{October 23, 2026}

\begin{document}
\begin{abstract}
Over any commutative ring, a module is finite projective if and only if
it is finitely presented and flat. We give the constructive proof using
finite relations, including the passage from one relation to a finite
system and the construction of a linear map through a surjection.
The definitions and constructions are those used in the formalisation.
\end{abstract}
\maketitle

\section{Definitions and statement}

Throughout, $R$ is a commutative ring with identity, and all maps between
modules are $R$-linear. The zero ring is allowed. All existence
assertions mean mere existence: a witness may be used in a proof of a
proposition, but its existence does not specify a chosen witness. A map
$q:P\to M$ is surjective if, for each $x\in M$, there merely exists
$v\in P$ such that $q(v)=x$.

For $n\geq0$, let $e_1,\ldots,e_n$ denote the usual basis of $R^n$.
We use the elementary identities
\[
 v=\sum_{i=1}^n v_i e_i,\qquad
 f\left(\sum_i a_i v_i\right)=\sum_i a_i f(v_i).
\]
In particular, a linear map from $R^n$ is determined by its values on
the $e_i$. These identities also hold for $n=0$, with empty sums.

\begin{definition}\label{def:fp}
A \emph{finite presentation} of $M$ consists of integers $n,m\geq0$,
a surjection $q:R^n\to M$, and vectors $a_1,\ldots,a_m\in R^n$ such that
\[
 q(a_j)=0\quad(1\leq j\leq m),\qquad
 q(v)=0\ \Longrightarrow\
 \exists c_1,\ldots,c_m\in R\quad v=\sum_j c_j a_j.
\]
The module $M$ is \emph{finitely presented} if such a presentation
merely exists.
\end{definition}

Writing $A:R^m\to R^n$ for the map whose columns are the $a_j$, this
is precisely a finite exact presentation
\[
 R^m\xrightarrow{A}R^n\xrightarrow{q}M\longrightarrow0.
\]
The containment $\operatorname{im}A\subseteq\ker q$ follows from
$q(a_j)=0$ and linearity; the other containment is supplied explicitly.

\begin{definition}\label{def:projective}
The module $M$ is \emph{finite projective} if there merely exist an
integer $k\geq0$ and maps
\[
 p:R^k\longrightarrow M,\qquad s:M\longrightarrow R^k
 \quad\text{with}\quad ps=\id_M.
\]
Thus finite projectivity is expressed as being a retract of a finite
free module.
\end{definition}

\begin{definition}\label{def:flat}
We use the single-relation criterion for \emph{flatness}: for every
$n\geq0$, every $q:R^n\to M$, and every $a\in R^n$ with $q(a)=0$,
there merely exist an integer $k\geq0$ and maps
\[
 B:R^n\longrightarrow R^k,\qquad u:R^k\longrightarrow M
 \quad\text{such that}\quad uB=q,\qquad B(a)=0.
\]
\end{definition}

To write this in terms of elements, put $x_i=q(e_i)$ and
$y_l=u(e_l)$. If the entries of $B$ are $b_{li}$, the criterion says
that a relation $\sum_i a_i x_i=0$ can be written as
\[
 x_i=\sum_l b_{li}y_l,\qquad
 \sum_i a_i b_{li}=0\quad\text{for every }l.
\]
Thus the given relation follows from finitely many relations among
scalars. Definition~\ref{def:flat} is the formulation of flatness used
in the formalisation.

\begin{theorem}\label{thm:main}
An $R$-module is finite projective if and only if it is finitely
presented and flat.
\end{theorem}

\section{One relation and finitely many relations}

\begin{lemma}\label{lem:systems}
Let $M$ be flat. Suppose $q:R^n\to M$ and
$a_1,\ldots,a_m\in R^n$ satisfy $q(a_j)=0$ for every $j$. Then there
merely exist $k\geq0$ and maps
\[
 B:R^n\longrightarrow R^k,\qquad u:R^k\longrightarrow M
\]
such that
\[
 uB=q,\qquad B(a_j)=0\quad(1\leq j\leq m).
\]
\end{lemma}
\begin{proof}
Induct on $m$. For $m=0$, take $k=n$, $B=\id_{R^n}$, and $u=q$.

For $m>0$, first apply flatness to $q(a_m)=0$. Obtain
\[
 B_0:R^n\to R^{k_0},\qquad u_0:R^{k_0}\to M,
 \qquad u_0B_0=q,\qquad B_0(a_m)=0.
\]
For $j<m$ we have $u_0(B_0(a_j))=q(a_j)=0$. Apply the induction
hypothesis to $u_0$ and these $m-1$ vectors. It gives
\[
 C:R^{k_0}\to R^k,\qquad u:R^k\to M,
 \qquad uC=u_0,\qquad C(B_0(a_j))=0\quad(j<m).
\]
Set $B=CB_0$. Then $uB=uCB_0=u_0B_0=q$. The first $m-1$ vectors
are killed by construction, and $B(a_m)=C(0)=0$.
\end{proof}

The finite system in this lemma is obtained by induction from the
single-relation definition; it is not an additional hypothesis on $M$.

\section{Factoring through a surjection}

\begin{lemma}\label{lem:factor}
Let $q:P\to M$ be surjective, and let $b:P\to N$ satisfy
\[
 q(v)=0\ \Longrightarrow\ b(v)=0.
\]
There is a unique linear map $s:M\to N$ such that $sq=b$.
\end{lemma}
\begin{proof}
If $q(v)=q(w)$, then $q(v-w)=0$, so $b(v-w)=0$ and $b(v)=b(w)$.

For each $x\in M$, there is therefore a unique $y\in N$ with the
property
\begin{equation}\label{eq:value}
 q(v)=x\ \Longrightarrow\ b(v)=y\qquad(v\in P).
\end{equation}
Existence follows by using one preimage of $x$; independence of that
preimage is the equality just proved. Uniqueness follows because $x$
has a preimage. Define $s(x)$ to be this unique value. This construction
uses unique choice of a value, rather than a choice of a preimage for
every $x$.

By \eqref{eq:value}, $s(q(v))=b(v)$. To prove additivity, use preimages
$q(v)=x$ and $q(w)=x'$ inside the proof of the equality. Then
\[
 s(x+x')=s(q(v+w))=b(v+w)=b(v)+b(w)=s(x)+s(x').
\]
Similarly, for $c\in R$,
\[
 s(cx)=s(q(cv))=b(cv)=c\,b(v)=c\,s(x).
\]
These are equalities in sets, so the mere existence of the finitely
many preimages suffices to prove them. Hence $s$ is linear.
Finally, any map satisfying $sq=b$ must take every $x$ to the unique
value in \eqref{eq:value}, proving uniqueness.
\end{proof}

\section{Finite presentation and flatness give a retract}

\begin{proof}
Suppose $M$ is finitely presented and flat. Since our conclusion is an
existence assertion, we may use a finite presentation as in
Definition~\ref{def:fp}. Lemma~\ref{lem:systems} applied to its relators
gives $k$, $B:R^n\to R^k$, and $u:R^k\to M$ with
\[
 uB=q,\qquad B(a_j)=0\quad(1\leq j\leq m).
\]
The map $B$ kills all of $\ker q$. Indeed, if $q(v)=0$, the presentation
gives $v=\sum_j c_j a_j$, and hence
\[
 B(v)=\sum_j c_j B(a_j)=0.
\]
By Lemma~\ref{lem:factor}, $B$ factors uniquely as
\[
 B=sq\qquad\text{for a linear }s:M\to R^k.
\]
Consequently $usq=uB=q$. For $x\in M$, use any preimage $q(v)=x$ to
obtain
\[
 us(x)=usq(v)=q(v)=x.
\]
Thus $us=\id_M$, and $(k,u,s)$ exhibits $M$ as a retract of $R^k$.
The required retract therefore merely exists.
\end{proof}

\section{A retract is finitely presented and flat}

\begin{proof}
Suppose $M$ is finite projective. We may use maps
\[
 p:R^k\to M,\qquad s:M\to R^k,\qquad ps=\id_M.
\]
We construct a finite presentation and verify the single-relation
criterion directly.

For the presentation put
\[
 h=\id_{R^k}-sp,\qquad a_i=h(e_i)=e_i-s(p(e_i))\quad(1\leq i\leq k).
\]
The map $p$ is surjective, with preimage $s(x)$ for $x\in M$.
Moreover,
\[
 p(a_i)=p(e_i)-psp(e_i)=0.
\]
If $p(v)=0$, then $h(v)=v-s(0)=v$. Writing $v=\sum_i v_i e_i$ and
using linearity gives
\[
 v=h(v)=\sum_i v_i h(e_i)=\sum_i v_i a_i.
\]
Thus the $a_i$ generate $\ker p$, and we have the finite presentation
\[
 R^k\xrightarrow{\,\id-sp\,}R^k\xrightarrow{p}M\longrightarrow0.
\]

To prove flatness, let $q:R^n\to M$ and let $a\in R^n$ satisfy
$q(a)=0$. Take
\[
 B=sq:R^n\to R^k,\qquad u=p:R^k\to M.
\]
Then $uB=psq=q$ and $B(a)=s(q(a))=s(0)=0$. This is exactly
Definition~\ref{def:flat}.

Both constructions are made from a retract. Mere existence of a
retract gives mere existence of the presentation, and, for each
relation separately, mere existence of its explanation. This proves
finite presentation and flatness, completing the theorem.
\end{proof}

\section*{Correspondence with the formalisation}

The proof above follows the four files below in
\code{sag/chatelet}. All definitions are over an arbitrary
commutative ring $R$ with a set as its underlying type.

\smallskip
\noindent
\begin{tabularx}{\textwidth}{@{}>{\raggedright\arraybackslash}p{0.43\textwidth}>{\raggedright\arraybackslash}X@{}}
\textbf{Mathematical statement} & \textbf{File and declaration} \\
\hline
Finite linear combinations and their images under linear maps &
\code{DescentModuleMaps.seg}: \code{Free.LinearMap.equation},
\code{Killed.theorem} \\[4pt]
Flatness, Definition~\ref{def:flat} &
\code{DescentFlat.seg}: \code{At.Flat} \\[4pt]
Finite systems, Lemma~\ref{lem:systems} &
\code{DescentFlat.seg}: \code{At.OneToMany.theorem} \\[4pt]
Factor through a surjection, Lemma~\ref{lem:factor} &
\code{DescentSurjectiveFactor.seg}: \code{At.factor},
\code{At.beta}, \code{At.linear} \\[4pt]
Finite presentation and finite projectivity &
\code{DescentProjective.seg}: \code{At.Presentation},
\code{At.Retract}, \code{At.FinitelyPresented},
\code{At.FiniteProjective} \\[4pt]
Finitely presented and flat implies finite projective &
\code{DescentProjective.seg}: \code{At.fp-flat-to-projective} \\[4pt]
A finite projective module is finitely presented and flat &
\code{DescentProjective.seg}: \code{At.projective-to-fp},
\code{At.projective-to-flat} \\[4pt]
Theorem~\ref{thm:main} &
\code{DescentProjective.seg}: \code{At.finite-projective-iff-fp-flat}
\end{tabularx}

\medskip
The main declaration is
\begin{center}
\code{DescentProjective.Over[R,setR].At[M].}\\
\code{finite-projective-iff-fp-flat}.
\end{center}
In \code{DescentFlat}, the data $k,B,u$ with $uB=q$ are recorded in
\code{FiniteFactor}; \code{Explanation} additionally records that $B$
kills the specified relators. In \code{DescentProjective}, the reverse
implication constructs the map $s$ of Section~4, and the forward
implication uses exactly the relators $(\id-sp)(e_i)$ of Section~5.

\end{document}
